We use a four-step proof spine centered on one canonical object.

\subsection*{Main proof spine}
\begin{enumerate}
  \item \textbf{Define the object.}
  Set
  \[
  \Phi(q):=\Gamma(q)(1-q)\zeta(q).
  \]

  \item \textbf{Define balance and correspondence target.}
  Balance means $\Phi(q)=0$.
  Target~A asks for correspondence between balanced strip points and nontrivial zeta zeros (\cref{thm:target-correspondence}).

  \item \textbf{State one exclusion target.}
  Target~B asks to prove
  \[
  \Phi(q)=0 \Longrightarrow \RePart(q)=\tfrac12
  \]
  for strip points (\cref{thm:target-uniqueness}).

  \item \textbf{Deduce RH conditionally.}
  Combine Targets~A and~B to obtain the conditional RH reduction (\cref{prop:rh-reduction-main-spine}).
\end{enumerate}

\subsection*{Status}
Step~1 is definitional.
Steps~2 and~3 are open theorem targets.
Therefore Step~4 is conditional, and RH is not claimed as proved.

\subsection*{Placement of secondary machinery}
Alternative operator/resolvent architectures are secondary.
If retained, they belong to supporting notes and not to this main proof spine.
